\documentclass[11pt, ukrainian]{article}
\setlength\textwidth{170mm} \setlength\hoffset{-25mm}
\setlength\textheight{237mm}
 \setlength\voffset{-18mm}%\setlength\textheight{277mm} \setlength\voffset{-38mm}
%\setlength\parindent{0mm}
\pagestyle{empty}
\usepackage[T2A]{fontenc}
\usepackage[utf8]{inputenc}
\usepackage[ukrainian]{babel}
\usepackage{graphicx}
\usepackage{caption}
\DeclareCaptionFormat{nocaption}{}
\captionsetup{format=nocaption,aboveskip=0pt,belowskip=0pt}
\usepackage[Export]{adjustbox} 
\adjustboxset{max size={0.9\linewidth}{0.9\paperheight}}
\begin{document}
%begin
{29.11-02.12.2022 Контрольна робота, варіант  \No \textbf { 1 }\par}
Увага! Відповіді мають бути розгорнутими і при цьому обидві вміщуватись на одну сторону аркушу формату А4.

{
\begin {large}

%beginitem
1. Що з наведеного переліку (якщо таке в переліку є):

\vspace{4mm}
\qquadпослідовність виконання інструкцій, прозорість за посиланнями, ледачі обчислення, анонімні функції

\vspace{4mm}
має найслабший логічний зв’язок з більшістю понять (відноситься до іншої парадигми та/або виривається з контексту чи не дуже сумісне)? Відповідь пояснити.

%enditem

%beginitem
2. Що є наслідком захисту м’ютексами занадто великого фрагменту коду?

%enditem
\end {large}}

\vspace{4mm}
%end



%begin
{29.11-02.12.2022 Контрольна робота, варіант  \No \textbf { 2 }\par}
Увага! Відповіді мають бути розгорнутими і при цьому обидві вміщуватись на одну сторону аркушу формату А4.

{
\begin {large}

%beginitem
1. Що з наведеного переліку (якщо таке в переліку є):

\vspace{4mm}
\qquadенергійні обчислення, виклик за необхідністю, прозорість за посиланнями, каррінг

\vspace{4mm}
має найслабший логічний зв’язок з більшістю понять (відноситься до іншої парадигми та/або виривається з контексту чи не дуже сумісне)? Відповідь пояснити.

%enditem

%beginitem
2. Що таке недетерміновано впорядковані (indeterminately sequenced) обчислення? Чи є їх наявність помилкою у програмі?  Відповідь пояснити.

%enditem
\end {large}}

\vspace{4mm}
%end



%begin
{29.11-02.12.2022 Контрольна робота, варіант  \No \textbf { 3 }\par}
Увага! Відповіді мають бути розгорнутими і при цьому обидві вміщуватись на одну сторону аркушу формату А4.

{
\begin {large}

%beginitem
1. Що з наступного не підтримує ані мова С++, ані її стандартна бібліотека: каррінг, виклик за необхідністю, замикання, ледачі обчислення, строгі (strict) функції? Відповідь пояснити.%enditem

%beginitem
2. Що є точкою входу в новий потік, що його запущено засобами бібліотеки thread? Чи може батьківський потік отримати з нього значення-результат? Чи можна і якщо можна, то як, у новий потік передати посилання на дані батьківського потоку?

%enditem
\end {large}}

\vspace{4mm}
%end



%begin
{29.11-02.12.2022 Контрольна робота, варіант  \No \textbf { 4 }\par}
Увага! Відповіді мають бути розгорнутими і при цьому обидві вміщуватись на одну сторону аркушу формату А4.

{
\begin {large}

%beginitem
1. Чи можна програмувати мовою Java у стилі структурного програмування? Відповідь пояснити.%enditem

%beginitem
2. Що з точки зору С++20 можна сказати про послідовність обчислень підвиразів у виразі f(R(13), new P(b+c), Q(a))) ? Вважати всі імена коректно означеними.

%enditem
\end {large}}

\vspace{4mm}
%end
\newpage
%begin
{29.11-02.12.2022 Контрольна робота, варіант  \No \textbf { 5 }\par}
Увага! Відповіді мають бути розгорнутими і при цьому обидві вміщуватись на одну сторону аркушу формату А4.

{
\begin {large}

%beginitem
1. Чи можна програмувати мовою Java у стилі декларативного програмування? Відповідь пояснити.%enditem

%beginitem
2. Чи завжди наслідком стану гонки (race condition) є невизначена поведінка програми? Відповідь пояснити.

%enditem
\end {large}}

\vspace{4mm}
%end



%begin
{29.11-02.12.2022 Контрольна робота, варіант  \No \textbf { 6 }\par}
Увага! Відповіді мають бути розгорнутими і при цьому обидві вміщуватись на одну сторону аркушу формату А4.

{
\begin {large}

%beginitem
1. Чи можна програмувати мовою С++ у стилі структурного програмування? Відповідь пояснити.%enditem

%beginitem
2. Які бібліотечні засоби С++20 та як здатні забезпечувати захист глобальних даних від стану гонки, якщо може виконуватись довільна кількість операцій доступу (читання та запису) до цих даних?

%enditem
\end {large}}

\vspace{4mm}
%end



%begin
{29.11-02.12.2022 Контрольна робота, варіант  \No \textbf { 7 }\par}
Увага! Відповіді мають бути розгорнутими і при цьому обидві вміщуватись на одну сторону аркушу формату А4.

{
\begin {large}

%beginitem
1. Що з наведеного переліку (якщо таке в переліку є):

\vspace{4mm}
\qquadфункції вищого порядку, каррінг, функції першого порядку, SECD-машина

\vspace{4mm}
має найслабший логічний зв’язок з більшістю понять (відноситься до іншої парадигми та/або виривається з контексту чи не дуже сумісне)? Відповідь пояснити.

%enditem

%beginitem
2. Які проблеми можуть виникнути за concurrent-доступу до даних, що зберігаються в структурі даних з операціями завантажити дані та ітератором, що забезпечує її обхід? Як їх можна уникнути?

%enditem
\end {large}}

\vspace{4mm}
%end



%begin
{29.11-02.12.2022 Контрольна робота, варіант  \No \textbf { 8 }\par}
Увага! Відповіді мають бути розгорнутими і при цьому обидві вміщуватись на одну сторону аркушу формату А4.

{
\begin {large}

%beginitem
1. Чи можна одну ту саму мову програмування використовувати для програмування у функціональному, імперативному та об’єктно-орієнтованому стилі одночасно? Відповідь пояснити.%enditem

%beginitem
2. У наслідок чого можуть виникати мертві блокування (deadlock)? Як з ними можна боротися? Які є бібліотечні засоби, що частково вирішують проблему?

%enditem
\end {large}}

\vspace{4mm}
%end
\newpage
%begin
{29.11-02.12.2022 Контрольна робота, варіант  \No \textbf { 9 }\par}
Увага! Відповіді мають бути розгорнутими і при цьому обидві вміщуватись на одну сторону аркушу формату А4.

{
\begin {large}

%beginitem
1. Чи можна програмувати мовою С++ у стилі декларативного програмування? Відповідь пояснити.%enditem

%beginitem
2. Які є засоби С++20 та прийоми, щоб забезпечити безпечне створення глобальних даних?

%enditem
\end {large}}

\vspace{4mm}
%end



%begin
{29.11-02.12.2022 Контрольна робота, варіант  \No \textbf { 10 }\par}
Увага! Відповіді мають бути розгорнутими і при цьому обидві вміщуватись на одну сторону аркушу формату А4.

{
\begin {large}

%beginitem
1. Чи підходить мова Haskell для програмування в імперативному стилі? Відповідь пояснити.%enditem

%beginitem
2. Як найкраще (з максимальним збереженням конкурентності) забезпечити захист глобальних даних від стану гонки, якщо над ними переважно виконуються операції читання і подекуди операції запису?

%enditem
\end {large}}

\vspace{4mm}
%end



%begin
{29.11-02.12.2022 Контрольна робота, варіант  \No \textbf { 11 }\par}
Увага! Відповіді мають бути розгорнутими і при цьому обидві вміщуватись на одну сторону аркушу формату А4.

{
\begin {large}

%beginitem
1. Що з наведеного переліку (якщо таке в переліку є):

\vspace{4mm}
\qquadнормальний порядок редукцій, замикання, каррінг, апплікативний порядок редукцій

\vspace{4mm}
має найслабший логічний зв’язок з більшістю понять (відноситься до іншої парадигми та/або виривається з контексту чи не дуже сумісне)? Відповідь пояснити.

%enditem

%beginitem
2. Які проблеми можуть виникнути за concurrent-доступу до даних, що зберігаються у двобічно зв’язаному списку з операціями вставити, додати та ітератором, що забезпечує його обхід? Як їх можна уникнути?

%enditem
\end {large}}

\vspace{4mm}
%end



%begin
{29.11-02.12.2022 Контрольна робота, варіант  \No \textbf { 12 }\par}
Увага! Відповіді мають бути розгорнутими і при цьому обидві вміщуватись на одну сторону аркушу формату А4.

{
\begin {large}

%beginitem
1. Що з наведеного переліку (якщо таке в переліку є):

\vspace{4mm}
\quadдекларативне програмування, енергійні обчислення, функції першого порядку, стан пам’яті

\vspace{4mm}
має найслабший логічний зв’язок з більшістю понять (відноситься до іншої парадигми та/або виривається з контексту чи не дуже сумісне)? Відповідь пояснити.

%enditem

%beginitem
2. Як співвідносяться потоки операційної системи з потоками-об’єктами thread? Який з методів join, detach є блокуючим?

%enditem
\end {large}}

\vspace{4mm}
%end
\newpage
%begin
{29.11-02.12.2022 Контрольна робота, варіант  \No \textbf { 13 }\par}
Увага! Відповіді мають бути розгорнутими і при цьому обидві вміщуватись на одну сторону аркушу формату А4.

{
\begin {large}

%beginitem
1. Чи можна програмувати мовою С++ у стилі декларативного програмування? Відповідь пояснити.%enditem

%beginitem
2. Чи мають відношення обіцянки (promise) до сопрограм? Відповідь пояснити.

%enditem
\end {large}}

\vspace{4mm}
%end



%begin
{29.11-02.12.2022 Контрольна робота, варіант  \No \textbf { 14 }\par}
Увага! Відповіді мають бути розгорнутими і при цьому обидві вміщуватись на одну сторону аркушу формату А4.

{
\begin {large}

%beginitem
1. Що з наведеного переліку (якщо таке в переліку є):

\vspace{4mm}
\qquadапплікативний порядок редукцій, строгі (strict) функції, каррінг, виклик за значенням

\vspace{4mm}
має найслабший логічний зв’язок з більшістю понять (відноситься до іншої парадигми та/або виривається з контексту чи не дуже сумісне)? Відповідь пояснити.

%enditem

%beginitem
2. Які бібліотечні засоби С++20 дозволяють не тільки запустити код в іншому потоці, а ще без особливих танців з бубнами дізнатись результат його виконання?

%enditem
\end {large}}

\vspace{4mm}
%end



%begin
{29.11-02.12.2022 Контрольна робота, варіант  \No \textbf { 15 }\par}
Увага! Відповіді мають бути розгорнутими і при цьому обидві вміщуватись на одну сторону аркушу формату А4.

{
\begin {large}

%beginitem
1. Що з наведеного переліку (якщо таке в переліку є):

\vspace{4mm}
\qquadнормальний порядок редукцій, ледачі обчислення, структурне програмування, анонімні функції

\vspace{4mm}
має найслабший логічний зв’язок з більшістю понять (відноситься до іншої парадигми та/або виривається з контексту чи не дуже сумісне)? Відповідь пояснити.

%enditem

%beginitem
2. Чи можна обмежити доступ потоку до глобальних змінних? А якщо дуже треба? Відповідь пояснити.

%enditem
\end {large}}

\vspace{4mm}
%end



%begin
{29.11-02.12.2022 Контрольна робота, варіант  \No \textbf { 16 }\par}
Увага! Відповіді мають бути розгорнутими і при цьому обидві вміщуватись на одну сторону аркушу формату А4.

{
\begin {large}

%beginitem
1. Чи можна програмувати мовою С++ у стилі структурного програмування? Відповідь пояснити.%enditem

%beginitem
2. Що з точки зору С++20 можна сказати про послідовність обчислень підвиразів у виразі f(a+b)+g(c+d) ? Вважати всі імена коректно означеними.

%enditem
\end {large}}

\vspace{4mm}
%end
\newpage
%begin
{29.11-02.12.2022 Контрольна робота, варіант  \No \textbf { 17 }\par}
Увага! Відповіді мають бути розгорнутими і при цьому обидві вміщуватись на одну сторону аркушу формату А4.

{
\begin {large}

%beginitem
1. Чи можна програмувати мовою Java у стилі декларативного програмування? Відповідь пояснити.%enditem

%beginitem
2. Фьючерси та обіцянки (promise): як за допомогою них можна організувати асинхронне виконання частин коду?

%enditem
\end {large}}

\vspace{4mm}
%end



%begin
{29.11-02.12.2022 Контрольна робота, варіант  \No \textbf { 18 }\par}
Увага! Відповіді мають бути розгорнутими і при цьому обидві вміщуватись на одну сторону аркушу формату А4.

{
\begin {large}

%beginitem
1. Чи можна одну ту саму мову програмування використовувати для програмування у функціональному, імперативному та об’єктно-орієнтованому стилі одночасно? Відповідь пояснити.%enditem

%beginitem
2. Які бібліотечні засоби С++20 та за якою схемою (кратко) здатні забезпечувати синхронізацію потоків?

%enditem
\end {large}}

\vspace{4mm}
%end



%begin
{29.11-02.12.2022 Контрольна робота, варіант  \No \textbf { 19 }\par}
Увага! Відповіді мають бути розгорнутими і при цьому обидві вміщуватись на одну сторону аркушу формату А4.

{
\begin {large}

%beginitem
1. Чи підходить мова Haskell для програмування в імперативному стилі? Відповідь пояснити.%enditem

%beginitem
2. Які бібліотечні засоби С++20 можна використати, що забезпечити передачу результатів обчислення від одного потоку до іншого? Яка схема їх використання?

%enditem
\end {large}}

\vspace{4mm}
%end



%begin
{29.11-02.12.2022 Контрольна робота, варіант  \No \textbf { 20 }\par}
Увага! Відповіді мають бути розгорнутими і при цьому обидві вміщуватись на одну сторону аркушу формату А4.

{
\begin {large}

%beginitem
1. Що з наведеного переліку (якщо таке в переліку є):

\vspace{4mm}
\qquadблок, стан пам’яті, побічний ефект, замикання

\vspace{4mm}
має найслабший логічний зв’язок з більшістю понять (відноситься до іншої парадигми та/або виривається з контексту чи не дуже сумісне)? Відповідь пояснити.

%enditem

%beginitem
2. Чи може потік мати доступ до автоматичної пам’яті іншого потоку? Відповідь пояснити.

%enditem
\end {large}}

\vspace{4mm}
%end
\newpage
%begin
{29.11-02.12.2022 Контрольна робота, варіант  \No \textbf { 21 }\par}
Увага! Відповіді мають бути розгорнутими і при цьому обидві вміщуватись на одну сторону аркушу формату А4.

{
\begin {large}

%beginitem
1. Чи можна програмувати мовою Java у стилі структурного програмування? Відповідь пояснити.%enditem

%beginitem
2. Як співвідносяться потоки операційної системи з потоками-об’єктами thread? Який з методів join, detach є блокуючим?

%enditem
\end {large}}

\vspace{4mm}
%end



%begin
{29.11-02.12.2022 Контрольна робота, варіант  \No \textbf { 22 }\par}
Увага! Відповіді мають бути розгорнутими і при цьому обидві вміщуватись на одну сторону аркушу формату А4.

{
\begin {large}

%beginitem
1. Що з наступного не підтримує ані мова С++, ані її стандартна бібліотека: каррінг, виклик за необхідністю, замикання, ледачі обчислення, строгі (strict) функції? Відповідь пояснити.%enditem

%beginitem
2. У наслідок чого можуть виникати мертві блокування (deadlock)? Як з ними можна боротися? Які є бібліотечні засоби, що частково вирішують проблему?

%enditem
\end {large}}

\vspace{4mm}
%end



%begin
{29.11-02.12.2022 Контрольна робота, варіант  \No \textbf { 23 }\par}
Увага! Відповіді мають бути розгорнутими і при цьому обидві вміщуватись на одну сторону аркушу формату А4.

{
\begin {large}

%beginitem
1. Чи підходить мова Haskell для програмування в імперативному стилі? Відповідь пояснити.%enditem

%beginitem
2. Які є засоби С++20 та прийоми, щоб забезпечити безпечне створення глобальних даних?

%enditem
\end {large}}

\vspace{4mm}
%end



%begin
{29.11-02.12.2022 Контрольна робота, варіант  \No \textbf { 24 }\par}
Увага! Відповіді мають бути розгорнутими і при цьому обидві вміщуватись на одну сторону аркушу формату А4.

{
\begin {large}

%beginitem
1. Чи можна програмувати мовою С++ у стилі декларативного програмування? Відповідь пояснити.%enditem

%beginitem
2. Які бібліотечні засоби С++20 дозволяють не тільки запустити код в іншому потоці, а ще без особливих танців з бубнами дізнатись результат його виконання?

%enditem
\end {large}}

\vspace{4mm}
%end
\newpage
%begin
{29.11-02.12.2022 Контрольна робота, варіант  \No \textbf { 25 }\par}
Увага! Відповіді мають бути розгорнутими і при цьому обидві вміщуватись на одну сторону аркушу формату А4.

{
\begin {large}

%beginitem
1. Чи можна програмувати мовою Java у стилі структурного програмування? Відповідь пояснити.%enditem

%beginitem
2. Чи може потік мати доступ до автоматичної пам’яті іншого потоку? Відповідь пояснити.

%enditem
\end {large}}

\vspace{4mm}
%end



%begin
{29.11-02.12.2022 Контрольна робота, варіант  \No \textbf { 26 }\par}
Увага! Відповіді мають бути розгорнутими і при цьому обидві вміщуватись на одну сторону аркушу формату А4.

{
\begin {large}

%beginitem
1. Що з наведеного переліку (якщо таке в переліку є):

\vspace{4mm}
\qquadнормальний порядок редукцій, SECD-машина, послідовність виконання інструкцій, каррінг

\vspace{4mm}
має найслабший логічний зв’язок з більшістю понять (відноситься до іншої парадигми та/або виривається з контексту чи не дуже сумісне)? Відповідь пояснити.

%enditem

%beginitem
2. Що є точкою входу в новий потік, що його запущено засобами бібліотеки thread? Чи може батьківський потік отримати з нього значення-результат? Чи можна і якщо можна, то як, у новий потік передати посилання на дані батьківського потоку?

%enditem
\end {large}}

\vspace{4mm}
%end



%begin
{29.11-02.12.2022 Контрольна робота, варіант  \No \textbf { 27 }\par}
Увага! Відповіді мають бути розгорнутими і при цьому обидві вміщуватись на одну сторону аркушу формату А4.

{
\begin {large}

%beginitem
1. Чи можна програмувати мовою С++ у стилі структурного програмування? Відповідь пояснити.%enditem

%beginitem
2. Які проблеми можуть виникнути за concurrent-доступу до даних, що зберігаються в структурі даних з операціями завантажити дані та ітератором, що забезпечує її обхід? Як їх можна уникнути?

%enditem
\end {large}}

\vspace{4mm}
%end



%begin
{29.11-02.12.2022 Контрольна робота, варіант  \No \textbf { 28 }\par}
Увага! Відповіді мають бути розгорнутими і при цьому обидві вміщуватись на одну сторону аркушу формату А4.

{
\begin {large}

%beginitem
1. Що з наступного не підтримує ані мова С++, ані її стандартна бібліотека: каррінг, виклик за необхідністю, замикання, ледачі обчислення, строгі (strict) функції? Відповідь пояснити.%enditem

%beginitem
2. Які бібліотечні засоби С++20 та як здатні забезпечувати захист глобальних даних від стану гонки, якщо може виконуватись довільна кількість операцій доступу (читання та запису) до цих даних?

%enditem
\end {large}}

\vspace{4mm}
%end
\newpage
%begin
{29.11-02.12.2022 Контрольна робота, варіант  \No \textbf { 29 }\par}
Увага! Відповіді мають бути розгорнутими і при цьому обидві вміщуватись на одну сторону аркушу формату А4.

{
\begin {large}

%beginitem
1. Що з наведеного переліку (якщо таке в переліку є):

\vspace{4mm}
\qquadфункціональне програмування, чисті функції, каррінг, рекурсія, прозорість за посиланнями

\vspace{4mm}
має найслабший логічний зв’язок з більшістю понять (відноситься до іншої парадигми та/або виривається з контексту чи не дуже сумісне)? Відповідь пояснити.

%enditem

%beginitem
2. Чи можна обмежити доступ потоку до глобальних змінних? А якщо дуже треба? Відповідь пояснити.

%enditem
\end {large}}

\vspace{4mm}
%end



%begin
{29.11-02.12.2022 Контрольна робота, варіант  \No \textbf { 30 }\par}
Увага! Відповіді мають бути розгорнутими і при цьому обидві вміщуватись на одну сторону аркушу формату А4.

{
\begin {large}

%beginitem
1. Що з наведеного переліку (якщо таке в переліку є):

\vspace{4mm}
\qquadфункції вищого порядку, анонімні функції, функції першого порядку, замикання

\vspace{4mm}
має найслабший логічний зв’язок з більшістю понять (відноситься до іншої парадигми та/або виривається з контексту чи не дуже сумісне)? Відповідь пояснити.

%enditem

%beginitem
2. Чи мають відношення обіцянки (promise) до сопрограм? Відповідь пояснити.

%enditem
\end {large}}

\vspace{4mm}
%end



%begin
{29.11-02.12.2022 Контрольна робота, варіант  \No \textbf { 31 }\par}
Увага! Відповіді мають бути розгорнутими і при цьому обидві вміщуватись на одну сторону аркушу формату А4.

{
\begin {large}

%beginitem
1. Що з наведеного переліку (якщо таке в переліку є):

\vspace{4mm}
\qquadвзаємодія об’єктів, чисті функції, абстракція даних, побічний ефект

\vspace{4mm}
має найслабший логічний зв’язок з більшістю понять (відноситься до іншої парадигми та/або виривається з контексту чи не дуже сумісне)? Відповідь пояснити.

%enditem

%beginitem
2. Які проблеми можуть виникнути за concurrent-доступу до даних, що зберігаються у двобічно зв’язаному списку з операціями вставити, додати та ітератором, що забезпечує його обхід? Як їх можна уникнути?

%enditem
\end {large}}

\vspace{4mm}
%end



%begin
{29.11-02.12.2022 Контрольна робота, варіант  \No \textbf { 32 }\par}
Увага! Відповіді мають бути розгорнутими і при цьому обидві вміщуватись на одну сторону аркушу формату А4.

{
\begin {large}

%beginitem
1. Чи можна програмувати мовою Java у стилі декларативного програмування? Відповідь пояснити.%enditem

%beginitem
2. Що з точки зору С++20 можна сказати про послідовність обчислень підвиразів у виразі f(a+b)+g(c+d) ? Вважати всі імена коректно означеними.

%enditem
\end {large}}

\vspace{4mm}
%end
\newpage
%begin
{29.11-02.12.2022 Контрольна робота, варіант  \No \textbf { 33 }\par}
Увага! Відповіді мають бути розгорнутими і при цьому обидві вміщуватись на одну сторону аркушу формату А4.

{
\begin {large}

%beginitem
1. Чи можна одну ту саму мову програмування використовувати для програмування у функціональному, імперативному та об’єктно-орієнтованому стилі одночасно? Відповідь пояснити.%enditem

%beginitem
2. Чи завжди наслідком стану гонки (race condition) є невизначена поведінка програми? Відповідь пояснити.

%enditem
\end {large}}

\vspace{4mm}
%end



%begin
{29.11-02.12.2022 Контрольна робота, варіант  \No \textbf { 34 }\par}
Увага! Відповіді мають бути розгорнутими і при цьому обидві вміщуватись на одну сторону аркушу формату А4.

{
\begin {large}

%beginitem
1. Чи можна програмувати мовою С++ у стилі структурного програмування? Відповідь пояснити.%enditem

%beginitem
2. Як найкраще (з максимальним збереженням конкурентності) забезпечити захист глобальних даних від стану гонки, якщо над ними переважно виконуються операції читання і подекуди операції запису?

%enditem
\end {large}}

\vspace{4mm}
%end



%begin
{29.11-02.12.2022 Контрольна робота, варіант  \No \textbf { 35 }\par}
Увага! Відповіді мають бути розгорнутими і при цьому обидві вміщуватись на одну сторону аркушу формату А4.

{
\begin {large}

%beginitem
1. Що з наведеного переліку (якщо таке в переліку є):

\vspace{4mm}
\qquadнормальний порядок редукцій, ледачі обчислення, структурне програмування, анонімні функції

\vspace{4mm}
має найслабший логічний зв’язок з більшістю понять (відноситься до іншої парадигми та/або виривається з контексту чи не дуже сумісне)? Відповідь пояснити.

%enditem

%beginitem
2. Які бібліотечні засоби С++20 можна використати, що забезпечити передачу результатів обчислення від одного потоку до іншого? Яка схема їх використання?

%enditem
\end {large}}

\vspace{4mm}
%end



%begin
{29.11-02.12.2022 Контрольна робота, варіант  \No \textbf { 36 }\par}
Увага! Відповіді мають бути розгорнутими і при цьому обидві вміщуватись на одну сторону аркушу формату А4.

{
\begin {large}

%beginitem
1. Чи можна програмувати мовою Java у стилі структурного програмування? Відповідь пояснити.%enditem

%beginitem
2. Що таке недетерміновано впорядковані (indeterminately sequenced) обчислення? Чи є їх наявність помилкою у програмі?  Відповідь пояснити.

%enditem
\end {large}}

\vspace{4mm}
%end
\newpage
\end{document}

